\section{Outer-class tree--double-theta obstruction}
\label{sec:outer-obstruction}

The separate tree--theta construction of \cite{KriebelTreeTheta2026} uses a
three-leaf level-2 theta blob with two reticulations and four switchings.  At
an exact algebraic point its core mixture factors as

\[
 M_{y,z}=P_{y+z}B_yB_z,
\]

so, after including the common pendant factors, the network tensor is a
three-leaf tree tensor.  All inverse-Fourier transition entries are strict.
The fixed normalized theta map has a nonzero \(15\)-by-\(15\) Jacobian
determinant

\[
 \frac{h(10h^2+1)}{2^{61}3^4 5^{14}}>0,
 \qquad h=5^{-1/4},
\]

whereas the tree map has rank \(9\).  The preimage theorem therefore makes
the theta-parameter locus mapping into the tree model a smooth
\(23\)-dimensional submanifold of codimension six.  A certified tangent and
the real-analytic implicit-function theorem perturb the collision into the
strict continuous-time cone.

Because the theta map is a submersion at the common point, its image contains
an ambient-open neighborhood.  A smaller tree output germ lies inside that
neighborhood and has a physical analytic theta section, proving the proper
directed containment in \cref{prop:outer}.

This example does not contradict \cref{thm:main}.  Exhaustive root
enumeration gives the theta mixed graph census \((7,0,7)\): it has no
tree-child admissible rooting and lies outside \(W_{\mathrm{TC}}\).  It is
therefore an outer obstruction, not the sharp transition from strong to weak
tree-childness.  It also explains why no proof here treats arbitrary
three-leaf nontrivial blobs as pointwise detectable.
